\begin{ouverture}[courbe]{8}{Fonctions, images et fonction carré}{Partie 2 · Proportionnalité et fonctions}
{Formules · Tableaux · Courbes · Images · Antécédents · $x\mapsto x^2$}
\objectif{Utiliser le vocabulaire et les notations des fonctions.}
\objectif{Passer d’une formule à un tableau ou à une courbe.}
\objectif{Représenter la fonction carré et résoudre graphiquement $x^2=a$.}
\prerequis{Repérage, substitution dans une expression et racine carrée.}
\end{ouverture}
\pagestyle{otMain}\essentiel\fontsize{11}{15}\selectfont
\essSection{1}{Une entrée et une image}
\begin{definition}Une fonction associe à chaque nombre autorisé une seule image. On écrit $f(x)$ pour l’image de $x$ par $f$, et $f:x\mapsto x^2+1$ pour décrire la fonction. Si $f(3)=8$, 8 est l’image de 3 et 3 est un antécédent de 8. Un nombre peut avoir plusieurs antécédents, ou aucun.\end{definition}
\begin{exemple}Pour $f(x)=x^2+1$, $f(-2)=5$. Les nombres $-2$ et 2 sont deux antécédents de 5. Aucun nombre n’a pour image 0, car $x^2+1\ge1$.\end{exemple}
\begin{savoirfaire}[Trois représentations]Une formule permet de calculer toute image ; un tableau de valeurs n’en donne que quelques-unes ; une courbe rassemble les points de coordonnées $(x\,;f(x))$. Un point appartient à la courbe lorsque son ordonnée est l’image de son abscisse.\end{savoirfaire}
\begin{attention}$f(2)$ ne signifie pas $f\times2$. Ne pas échanger l’image et l’antécédent dans une phrase. Un tableau partiel ne permet pas d’inventer l’image d’une entrée non indiquée.\end{attention}
\begin{propriete}Comparer deux fonctions graphiquement revient à comparer leurs ordonnées pour une même abscisse. Les intersections correspondent aux entrées donnant la même image.\end{propriete}
\clearpage
\essSection{2}{Lire une courbe et la fonction carré}
\begin{savoirfaire}[Lire une image ou des antécédents]Partir d’une abscisse, rejoindre la courbe verticalement, puis lire l’ordonnée. Pour chercher les antécédents d’un nombre, partir de cette ordonnée et repérer toutes les intersections horizontales avec la courbe. Une lecture graphique est souvent approchée.\end{savoirfaire}
\begin{definition}La fonction carré est la fonction $x\mapsto x^2$. Sa courbe, symétrique par rapport à l’axe des ordonnées, est une parabole qui passe par l’origine et reste au-dessus de l’axe des abscisses.\end{definition}
\begin{propriete}Résoudre graphiquement $x^2=a$, c’est chercher les abscisses des points d’intersection de la parabole avec la droite horizontale d’équation $y=a$ : deux points si $a>0$, un seul si $a=0$, aucun si $a<0$.\end{propriete}
\begin{savoirfaire}[Modéliser]Nommer la variable et son unité, exprimer la quantité cherchée, puis préciser les entrées possibles. Un nombre de billets est entier ; une longueur doit rester positive. Pour une cuve de base 12 m$^2$, le volume est $V(h)=12h$, pour $h$ compris entre 0 et la hauteur maximale.\end{savoirfaire}
\clearpage\rubrique{Découverte}
\begin{blocLivre}{ex-8-1}\begin{exercice}[1]{Calculer des images}
Pour $f:x\mapsto3x-4$, calculer $f(0)$, $f(2)$ et $f(-3)$.
\end{exercice}\end{blocLivre}
\begin{blocLivre}{ex-8-2}\begin{exercice}[1]{Lire un tableau}
À l’aide du tableau, donner $f(-2)$ et les antécédents de 1 qui y figurent.\begin{center}\renewcommand{\arraystretch}{1.35}\begin{tabular}{c|cccc}\toprule
$x$ & $-2$ & $0$ & $1$ & $3$\\\midrule
$f(x)$ & $5$ & $1$ & $-1$ & $1$\\
\bottomrule\end{tabular}\end{center}
\end{exercice}\end{blocLivre}
\begin{blocLivre}{ex-8-3}\begin{exercice}[1]{Traduire une phrase}
Écrire avec une égalité : « 7 est l’image de $-1$ par $g$ ». Dire ensuite ce que signifie $g(4)=7$. Ces deux informations sont-elles compatibles ?
\end{exercice}\end{blocLivre}
\begin{blocLivre}{ex-8-4}\begin{exercice}[1]{Un point sur une courbe}
Pour $f(x)=x^2-2$, les points $A(3\,;7)$ et $B(-2\,;1)$ appartiennent-ils à la courbe ?
\end{exercice}\end{blocLivre}
\begin{blocLivre}{ex-8-5}\begin{exercice}[1]{Un antécédent algébrique}
Pour $h(x)=5x+2$, déterminer l’antécédent de 17.
\end{exercice}\end{blocLivre}

\clearpage\rubrique{Entraînement}
\begin{blocLivre}{ex-8-6}\begin{exercice}[2]{Une formule et un programme}
On choisit un nombre, on ajoute 3, puis on élève le résultat au carré. Donner la fonction obtenue et les images de 1 et $-3$.
\end{exercice}\end{blocLivre}
\begin{blocLivre}{ex-8-7}\begin{exercice}[2]{Lire une ligne brisée}
La courbe est formée des segments dessinés, avec les sommets sur les graduations indiquées. Lire $f(-1)$ et tous les antécédents de 1.
\label{fig-ex-8-7}\figureGeo{fonction-brisee}
\end{exercice}\end{blocLivre}
\begin{blocLivre}{ex-8-8}\begin{exercice}[2]{Plusieurs ou aucun antécédent}
Pour $f(x)=x^2$, déterminer les antécédents de 16, de 0 et de $-1$.
\label{fig-ex-8-8}\figureGeo{fonction-carre}
\end{exercice}\end{blocLivre}
\begin{blocLivre}{ex-8-9}\begin{exercice}[2]{Comparer deux représentations}
Les fonctions $f(x)=x+2$ et $g(x)=2x-1$ sont représentées. Pour quelle valeur de $x$ leurs images sont-elles égales ? Vérifier par un calcul.
\label{fig-ex-8-9}\figureGeo{fonction-intersection}
\end{exercice}\end{blocLivre}
\begin{blocLivre}{ex-8-10}\begin{exercice}[2]{Représenter la fonction carré}
Compléter un tableau de valeurs de $x\mapsto x^2$ pour $x$ égal à $-3$, $-2$, $-1$, $-0{,}5$, 0, 0,5, 1, 2 et 3, puis tracer la courbe dans le repère.
\label{fig-ex-8-10}\figureGeo{carre-trace}
\end{exercice}\end{blocLivre}

\clearpage\rubrique{Problèmes}
\begin{blocLivre}{ex-8-11}\begin{exercice}[3]{Un réservoir}
Un réservoir cylindrique de base d’aire 0,8 m$^2$ a une hauteur maximale de 1,5 m. Exprimer son volume en fonction de la hauteur $h$, puis trouver la hauteur correspondant à 600 L.
\label{fig-ex-8-11}\figureGeo{cuve-hauteur}
\end{exercice}\end{blocLivre}
\begin{blocLivre}{ex-8-12}\begin{exercice}[3]{Résoudre graphiquement $x^2=a$}
À l’aide de la courbe de la fonction carré, lire des valeurs approchées des solutions de $x^2=5$, puis donner leurs valeurs exactes et leurs arrondis au dixième. Combien de solutions ont $x^2=0$ et $x^2=-2$ ?
\label{fig-ex-8-12}\figureGeo{carre-equation}
\end{exercice}\end{blocLivre}
\begin{blocLivre}{ex-8-13}\begin{exercice}[3]{Des aires possibles}
Un rectangle de périmètre 40 cm a un côté de longueur $x$. Exprimer son aire. Quelles sont les entrées possibles ? Donner les images de 3 et de 17.
\label{fig-ex-8-13}\figureGeo{rectangle-fonction}
\end{exercice}\end{blocLivre}
\begin{blocLivre}{ex-8-14}\begin{exercice}[3]{Une lecture de déplacement}
La courbe donne la distance au point de départ pendant une promenade. Les segments ont pour sommets $(0\,;0)$, $(10\,;600)$, $(15\,;600)$ et $(25\,;0)$, en minutes et mètres. Décrire les phases et calculer la distance totale parcourue sur cet aller-retour rectiligne.
\label{fig-ex-8-14}\figureGeo{promenade}
\end{exercice}\end{blocLivre}
\begin{blocLivre}{ex-8-15}\begin{exercice}[3]{Résoudre par un produit nul}
Pour $f(x)=(x-2)(3x+6)$, déterminer les antécédents de 0. Donner ensuite une image strictement négative.
\end{exercice}\end{blocLivre}

\clearpage
